\documentclass[12pt,letterpaper]{article}
\usepackage[utf8]{inputenc}
\usepackage[spanish]{babel}
\usepackage{graphicx}
\usepackage{amsmath}
\usepackage{geometry}
\usepackage{url}
\geometry{top=2.5cm, bottom=2.5cm, left=2.5cm, right=2.5cm}

\title{\textbf{Visualización y análisis de funciones de dos variables: dominio, curvas de nivel, trazas y tasas de cambio usando MATLAB}}
\author{Nora Carolina Servellón Palma}
\date{Abril 2026}

\begin{document}

\maketitle


\vspace{0.5cm}

Este proyecto consiste en implementar herramientas computacionales en MATLAB para explorar conceptos fundamentales del cálculo de funciones de dos variables, tomando como base el capítulo 15.1 del libro de Cálculo de varias variables. Se generarán visualizaciones de:
\begin{enumerate}
    \item Dominio de una función
    \item Gráficas y trazas de superficies (paraboloide elíptico, paraboloide hiperbólico, plano)
    \item Mapas de contorno y espaciamiento
Tasa de cambio promedio y gradiente numérico
\item Tasa de cambio promedio y gradiente numérico
\end{enumerate}

Se incluye también un ejemplo de función tabular (densidad del agua de mar) para mostrar que no todas las funciones tienen expresión analítica. El producto esperado es un reporte técnico en LaTeX con figuras propias generadas en MATLAB, el código fuente disponible en GitHub y una presentación en Beamer. El proyecto demuestra cómo la computación gráfica facilita la comprensión de conceptos abstractos del cálculo multivariable.

\vspace{1cm}

\section*{Palabras clave}
Funciones de dos variables, curvas de nivel, gradiente, MATLAB, visualización científica.

\begin{thebibliography}{9}
\bibitem{stewart}
  J. Stewart, \emph{Cálculo de varias variables}, 7ª ed., Cengage Learning, 2012, Capítulo 15.1.
\bibitem{matlab}
  MathWorks, \emph{MATLAB Documentation}, disponible en \url{https://www.mathworks.com/MATLAB Drive/untitled.m}.
\end{thebibliography}


\end{document}
